\documentclass{MMCStyle}
\graphicspath{{../figures/}}

\bianhao{202601139}
\tihao{A题}
\timu{长江十年禁渔下鱼类食物链恢复、混沌阈值与入侵风险的耦合建模}
\keyword{长江十年禁渔；四大家鱼；长江江豚；长江鲟；Hastings-Powell模型；李雅普诺夫指数}
\bibliographystyle{gbt7714-numerical}

\begin{document}
\begin{abstract}
围绕长江十年禁渔背景下鱼类生态系统的恢复与失衡问题，本文构建了“资源--功能群模块、珍稀物种模块和三级食物链混沌模块”三层框架。底层采用四类资源与四大家鱼功能群的耦合常微分方程，描述禁渔后资源变化向鱼群的传导；在此基础上，以食源约束、人工放流和通道阻隔共同驱动的模型分析长江江豚与长江鲟的恢复。食物链动力学部分采用 Hastings-Powell 三级模型，考察稳定、周期振荡及更复杂的不规则长期行为，并在问题四中以最大李雅普诺夫指数识别混沌候选区间。工业污染和外来入侵物种则被纳入情景扩展。问题一进一步选取四大家鱼总量与单网次捕获量代理（CPUE*）作为可观测指标，其中 2025 年“四大家鱼恢复到禁渔前约 1.8 倍”作为硬校准锚点，湖北段“单网次捕获量提升 120\%”仅用于拟合外的一致性检验。单区校准后，四大家鱼总量在 2025 年达到基线的 1.805 倍，CPUE* 达到 1.838 倍，对应相对误差分别为 0.26\% 和 16.47\%；与此同时，模型中的水草末值降至 0.204 倍，提示在当前参数组下恢复伴随底层承压上升。基线修复情景下，江豚和长江鲟末期规模分别达到初值的 1.133 倍和 1.140 倍。污染与入侵叠加后，内部综合健康指数由 0.686 降至 0.575；权重扰动后，前两类情景排序仍存在交换，提示该指标并非完全稳健。据此，本文建议将“禁渔”与“水质治理、通道修复、定向增殖放流和入侵物种压制”联动实施。
\end{abstract}

\newpage
\tableofcontents
\newpage
\pagenumbering{arabic}
\pagestyle{fancy}
